\begin{afterword}
\summaryhead

The post begins with the author's first break with Judaism: as a small child, the author was told a Hebrew prayer works without understanding the words. It then gives cases of purposes lost along the way. Students study what they will forget to get a degree. Bay Area teachers spend little time on science, the author says because of the tests required by the No Child Left Behind Act. Soviet shoe factories met quotas with tiny shoes. Medical researchers aim at the p<0.05 needed for publication.

A single person, the author argues, keeps track of why each step of a plan matters. Organizations can reward only what is measurable now, so each link in a chain of incentives pursues its own measure, and the final purpose is lost; the Act's chain from politicians to children is traced as an example. Quoting Musashi, the post asks readers to see how every rule and subgoal serves the final goal, and ends with the author's own strong dislike of lost purposes, which the author hopes can be taught.

\responsehead

The post's main idea is sound and old. When an organization rewards a measurable step, people pursue the step and the goal behind it is lost. Donald Campbell stated this in 1976 about exactly the post's main case, test scores in schools, and the idea is known as Campbell's law. The post does not mention it. What the post adds is the contrast with a single mind, which follows its own plans without losing the goal.

The cases are uneven. The Bay Area figures are reported correctly, and the newspaper report of the survey links the teachers who taught no science to schools that had missed the Act's reading and math targets. But the next sentence, that ``Virtually all classroom time is now spent on preparing for tests,'' has no source and is overstated: a national survey the same year found elementary schools still spending about three hours a week on science. It also found the milder trend the post describes, with 44\% of districts cutting other subjects to add reading and mathematics. The claim that school systems ``seem'' to be collapsing is offered without evidence, and the national assessment of 2007 recorded fourth-grade scores higher than in any earlier year. The medical case is stated with the hedge its source warrants. The physics joke rests on a threshold, p<0.0001, that I could not find; the known convention is for claims of discovery in particle physics, and is far stricter. The Soviet stories have no source I could reach.

The account of each link in the chain of incentives, from politicians to textbook committees to teachers, is given as fact. Some of it may be right, but none of it is sourced, and its most specific claims, such as teachers not getting through ``a fourth of the textbook,'' come with no evidence.

The contrast between the author and other people rests on self-report. The author ``reflexively'' asks where a judgment comes from, and has an ``exceptionally strong abhorrence to lost purposes''; the evidence that others care less is that the author's classmates in kindergarten or first grade said their prayers.

In short: an old and sound idea, stated without its history, and supported by cases of mixed quality, of which the strongest claims, near-total test preparation and collapsing schools, have no evidence.
\end{afterword}
